% ============================================================================
% WORKBOOK — Section 1.2: Recognizing Proportional Relationships
% CCSS 7.RP.A.2.A
% Practice-focused reinforcement for the study guide topic
% ============================================================================
\section{Recognizing Proportional Relationships}
\topicTitle{Recognizing Proportional Relationships}

\begin{quickReview}{Recognizing Proportional Relationships}
Two quantities are \textbf{proportional} if every pair shares the \textbf{same ratio}.

\begin{itemize}[leftmargin=*]
    \item \textbf{Table test:} Divide $y \div x$ for every row. If you get the same quotient every time, it is proportional.
    \item \textbf{Graph test:} The points form a \textbf{straight line through the origin} $(0,0)$.
\end{itemize}

\medskip
\textbf{Example:} $(2, 6),\; (3, 9),\; (5, 15)$. Ratios: $\frac{6}{2}=3$, $\frac{9}{3}=3$, $\frac{15}{5}=3$ — proportional!

\smallskip
\textbf{Key idea:} Even one mismatched ratio means the relationship is \textbf{not proportional}.
\end{quickReview}

% ── Warm-Up ──────────────────────────────────────────────────────────────────
\begin{practiceBox}[funGreen]{\faSun~Warm-Up}

\practiceHeader[funGreen]{Find the ratio $y \div x$ for each pair.}

\begin{multicols}{2}
\prob $x = 4$, $y = 12$ \answerBlank[2.5cm]
\answerExplain{$3$}{Divide $y$ by $x$: $12 \div 4 = 3$.}

\prob $x = 5$, $y = 20$ \answerBlank[2.5cm]
\answerExplain{$4$}{Divide $y$ by $x$: $20 \div 5 = 4$.}

\prob $x = 3$, $y = 7.5$ \answerBlank[2.5cm]
\answerExplain{$2.5$}{Divide $y$ by $x$: $7.5 \div 3 = 2.5$.}

\prob $x = 6$, $y = 18$ \answerBlank[2.5cm]
\answerExplain{$3$}{Divide $y$ by $x$: $18 \div 6 = 3$.}

\prob $x = 10$, $y = 25$ \answerBlank[2.5cm]
\answerExplain{$2.5$}{Divide $y$ by $x$: $25 \div 10 = 2.5$.}

\prob $x = 8$, $y = 32$ \answerBlank[2.5cm]
\answerExplain{$4$}{Divide $y$ by $x$: $32 \div 8 = 4$.}
\end{multicols}
\end{practiceBox}

% ── Practice A: Table Test ───────────────────────────────────────────────────
\begin{practiceBox}{Proportional or Not? (Tables)}

\practiceHeader{Use the table test: divide $y \div x$ for each row. Write ``proportional'' or ``not proportional'' and explain.}

\prob \begin{center}
\begin{tabular}{|c|c|}
\hline $x$ & $y$ \\ \hline $2$ & $8$ \\ $5$ & $20$ \\ $9$ & $36$ \\ \hline
\end{tabular}
\end{center}
\answerBlank[5cm]
\answerExplain{Proportional}{Ratios: $8 \div 2 = 4$, $20 \div 5 = 4$, $36 \div 9 = 4$. Every ratio equals $4$, so the relationship is proportional.}

\prob \begin{center}
\begin{tabular}{|c|c|}
\hline $x$ & $y$ \\ \hline $3$ & $9$ \\ $4$ & $11$ \\ $6$ & $15$ \\ \hline
\end{tabular}
\end{center}
\answerBlank[5cm]
\answerExplain{Not proportional}{Ratios: $9 \div 3 = 3$, $11 \div 4 = 2.75$, $15 \div 6 = 2.5$. The ratios are not equal, so it is not proportional.}

\prob \begin{center}
\begin{tabular}{|c|c|}
\hline $x$ & $y$ \\ \hline $4$ & $6$ \\ $8$ & $12$ \\ $12$ & $18$ \\ \hline
\end{tabular}
\end{center}
\answerBlank[5cm]
\answerExplain{Proportional}{Ratios: $6 \div 4 = 1.5$, $12 \div 8 = 1.5$, $18 \div 12 = 1.5$. Every ratio equals $1.5$, so the relationship is proportional.}

\prob \begin{center}
\begin{tabular}{|c|c|}
\hline $x$ & $y$ \\ \hline $1$ & $5$ \\ $2$ & $8$ \\ $3$ & $11$ \\ \hline
\end{tabular}
\end{center}
\answerBlank[5cm]
\answerExplain{Not proportional}{Ratios: $5 \div 1 = 5$, $8 \div 2 = 4$, $11 \div 3 \approx 3.67$. The ratios differ. The pattern adds $3$ each time (constant change) but that does not make a constant ratio.}

\prob \begin{center}
\begin{tabular}{|c|c|}
\hline $x$ & $y$ \\ \hline $\frac{1}{2}$ & $3$ \\ $1$ & $6$ \\ $2$ & $12$ \\ \hline
\end{tabular}
\end{center}
\answerBlank[5cm]
\answerExplain{Proportional}{Ratios: $3 \div \frac{1}{2} = 6$, $6 \div 1 = 6$, $12 \div 2 = 6$. Every ratio equals $6$, so it is proportional.}

\end{practiceBox}

% ── Practice B: True or False ────────────────────────────────────────────────
\begin{practiceBox}[funPurple]{True or False?}

\trueOrFalse{If $y \div x$ gives different values for different rows in a table, the relationship could still be proportional.}
\answerExplain{False}{A proportional relationship requires the \emph{same} ratio $y \div x$ for every pair. Even one different ratio means it is not proportional.}

\trueOrFalse{A straight line through the origin $(0, 0)$ always represents a proportional relationship.}
\answerExplain{True}{A proportional relationship graphs as a straight line that passes through the origin. Both conditions (straight line and through the origin) are needed, and both are satisfied here.}

\trueOrFalse{The equation $y = 3x + 2$ represents a proportional relationship.}
\answerExplain{False}{A proportional equation has the form $y = kx$ with no added constant. The $+2$ means the line does not pass through $(0, 0)$, so it is not proportional.}

\trueOrFalse{If a graph is a straight line that crosses the $y$-axis at $4$, the relationship is proportional.}
\answerExplain{False}{A proportional graph must pass through the origin $(0, 0)$. A $y$-intercept of $4$ means it does not start at zero, so the relationship is not proportional.}

\end{practiceBox}

% ── Practice C: Real-World Scenarios ─────────────────────────────────────────
\begin{practiceBox}[funOrange]{\faGlobe~Real-World Scenarios}

\practiceHeader[funOrange]{Decide if each situation is proportional. Explain your reasoning.}

\prob A gym charges $\$25$ per month with no sign-up fee. Is monthly cost proportional to the number of months? \answerBlank[5cm]
\answerExplain{Yes}{With no sign-up fee, cost $= 25 \times$ months. Every month costs the same per-month amount, and $0$ months costs $\$0$. The ratio is always $25$, so it is proportional.}

\prob A taxi ride costs $\$3$ flat fee plus $\$2$ per mile. Is total fare proportional to miles? \answerBlank[5cm]
\answerExplain{No}{The $\$3$ flat fee means a $0$-mile ride costs $\$3$, not $\$0$. The ratio changes with each mile (e.g., $5/1 = 5$, $7/2 = 3.5$), so it is not proportional.}

\prob A store sells pens at $\$0.50$ each. Is cost proportional to the number of pens? \answerBlank[5cm]
\answerExplain{Yes}{Each pen costs the same amount, there is no extra fee, and $0$ pens cost $\$0$. The ratio is always $0.50$, so cost is proportional to the number of pens.}

\prob A phone plan charges $\$10$/month plus $\$0.03$ per text. Is the monthly bill proportional to the number of texts? \answerBlank[5cm]
\answerExplain{No}{The $\$10$ base fee does not depend on texts. At $0$ texts the bill is $\$10$, not $\$0$. The per-text ratio changes, so the relationship is not proportional.}

\end{practiceBox}

% ── Challenge ────────────────────────────────────────────────────────────────
\begin{challengeBox}
\prob A table shows $(2, 6)$, $(4, 12)$, $(5, \; ?)$, $(8, 24)$. If the relationship is proportional, find the missing $y$-value. If it is not proportional, explain why.
\answerBlank[4cm]
\answerExplain{$y = 15$}{Check the ratio: $6 \div 2 = 3$, $12 \div 4 = 3$, $24 \div 8 = 3$. All ratios equal $3$, so it is proportional with $k = 3$. The missing value is $y = 3 \times 5 = 15$.}

\prob Create your own table of four $(x, y)$ pairs that is \textbf{not proportional} but where the first two pairs have the same ratio $y \div x = 2$. Explain why the whole table is not proportional.
\answerExplain{Answers vary; e.g., $(1,2), (3,6), (4,9), (5,11)$}{The first two rows give $2 \div 1 = 2$ and $6 \div 3 = 2$, but the third or fourth row has a different ratio (e.g., $9 \div 4 = 2.25$). One mismatched ratio is enough to make the whole relationship not proportional.}
\end{challengeBox}

\encouragement{You can spot a proportional relationship from a mile away now --- great detective work!}
